% Commutative diagram from the long-exact-sequence chapter.
% Source: book/tex/homology/long-exact.tex (lines 721-727).
\documentclass[border=2mm]{standalone}
\usepackage{amssymb,amsmath}
\usepackage{tikz-cd}
\newcommand{\ZZ}{\mathbb{Z}}
\newcommand{\eps}{\varepsilon}
\newcommand{\id}{\mathrm{id}}
\begin{document}
	\begin{tikzcd}
		& 0 \ar[d] & 0 \ar[d] & 0 \ar[d] \\
		0 \ar[r] & A_1 \ar[r] \ar[d] & B_1 \ar[r] \ar[d] & C_1 \ar[r] \ar[d] & 0 \\
		0 \ar[r] & A_2 \ar[r] \ar[d] & B_2 \ar[r] \ar[d] & C_2 \ar[r] \ar[d] & 0 \\
		0 \ar[r] & A_3 \ar[r] \ar[d] & B_3 \ar[r] \ar[d] & C_3 \ar[r] \ar[d] & 0 \\
		& 0 & 0 & 0 &
	\end{tikzcd}
\end{document}
